Fix the grammar and English of the given text using the guidance here.

* Use UK English, e.g. analyse rather than analyze.
* Include a ~ between citations and references and the text. e.g. Fiat--Shamir~\cite{}
* Include full words such as Figure instead of Fig. Section instead of Sec or \S
* Quotations should be as: ``I'm in a quote''.
* Use sentence casing for paragraphs, e.g., /paragraph{Hello world} should be /paragraph{Hello World.}
* Use frontend rather than front-end
* Use backend rather than back-end
* et al. is to be replaced by \etal. e.g., Ma et al. \cite{x} should be Ma~\etal~\cite{x}
* The hyphen (-) is the shortest mark, used to join compound words like "mother-in-law" or break words across lines. The en dash (–) is slightly longer, representing ranges or connections between words, as in "pages 45–52" or "New York–London flight." The em dash (—) is the longest dash, used to create a strong break in sentence structure—like this—for emphasis, interruptions, or abrupt changes in thought. In typing, you're using the right approach: two hyphens (--) for an en dash and three hyphens (---) for an em dash, which many word processors will automatically convert to proper dashes, though dedicated keyboard shortcuts also exist for each. Use three hypthens for em dash, two en dash and etc.

Outout in raw latex.